\chapter{Legislação Específica do TCE-MA}

\section{Sistema normativo do controle externo maranhense}

A atuação do Tribunal de Contas do Estado do Maranhão (TCE-MA) é disciplinada por um conjunto hierarquizado de normas. A Constituição Federal fornece o modelo geral dos arts. 70 a 75; a Constituição do Estado do Maranhão estrutura o controle externo estadual; a Lei Estadual nº 8.258/2005 estabelece a Lei Orgânica do Tribunal; o Regimento Interno disciplina organização e procedimento; e leis e instruções normativas específicas detalham a estrutura administrativa e determinados processos de fiscalização.

Essa hierarquia é juridicamente relevante. O Regimento Interno não pode ampliar competência constitucional nem contrariar a Lei Orgânica. Da mesma forma, uma instrução normativa pode disciplinar procedimentos dentro da competência do Tribunal, mas deve respeitar as normas superiores.

\begin{teoria}[estrutura normativa]
\[
\text{CF/1988}\rightarrow
\text{Constituição do MA}\rightarrow
\text{Lei Orgânica do TCE-MA}\rightarrow
\text{Regimento Interno}\rightarrow
\text{atos normativos específicos}
\]
A solução de uma questão exige identificar \textbf{qual nível normativo} disciplina o tema e se existe alteração posterior da redação originalmente estudada.
\end{teoria}

\section{Posição constitucional do TCE-MA}

O TCE-MA integra o sistema de controle externo do Estado. A Assembleia Legislativa exerce o controle externo estadual com auxílio do Tribunal, sem relação de subordinação hierárquica.

A expressão \emph{auxílio} deve ser compreendida em sentido constitucional: o Tribunal exerce competências próprias, como fiscalização, julgamento de contas de responsáveis, apreciação de atos sujeitos a registro, aplicação de sanções e realização de auditorias, nos termos da Constituição e da legislação estadual.

O TCE-MA não integra o Poder Judiciário. Suas decisões são administrativas e podem produzir efeitos jurídicos próprios, mas permanecem sujeitas a controle jurisdicional quanto à juridicidade.

\section{Lei Orgânica do TCE-MA — Lei Estadual nº 8.258/2005}

A Lei Orgânica constitui o núcleo legal do funcionamento do Tribunal. Ela disciplina jurisdição, competências, fiscalização, processos, decisões, sanções, recursos e organização institucional.

\subsection{Jurisdição de contas}

A jurisdição do Tribunal alcança as pessoas e situações definidas pela Constituição e pela Lei Orgânica. O elemento central é a relação com recursos, bens e valores públicos: quem arrecada, guarda, gerencia, administra ou utiliza recursos públicos pode sujeitar-se ao dever de prestar contas e à responsabilização perante o Tribunal nas hipóteses legais.

A jurisdição não se restringe a servidores efetivos. Particulares, entidades privadas e outros agentes podem estar sujeitos ao controle quando administram recursos públicos ou assumem obrigações relacionadas ao erário.

\subsection{Competências de fiscalização}

A Lei Orgânica concretiza competências de fiscalização sobre, entre outros objetos:
\begin{itemize}
\item contas de governo e contas de gestão;
\item receitas e despesas;
\item licitações e contratos;
\item atos de pessoal;
\item transferências e convênios;
\item patrimônio;
\item renúncia de receitas;
\item sistemas e controles administrativos;
\item aplicação de recursos estaduais e municipais submetidos à jurisdição do Tribunal.
\end{itemize}

\section{Contas de governo e contas de gestão}

A distinção entre contas de governo e contas de gestão é fundamental.

\termo{Contas de governo} retratam a execução global da política governamental e das finanças do chefe do Poder Executivo. O Tribunal emite parecer prévio, que subsidia o julgamento político pelo Poder Legislativo competente.

\termo{Contas de gestão} estão relacionadas à administração direta de recursos por ordenadores, gestores e demais responsáveis. Essas contas são submetidas ao julgamento do Tribunal nas hipóteses legais.

\begin{atencao}[efeitos distintos]
Parecer prévio e julgamento de contas são atos diferentes. A existência de parecer técnico do Tribunal sobre contas de governo não transforma o Tribunal no órgão político que realiza o julgamento final dessas contas.
\end{atencao}

\section{Fiscalizações e instrumentos de controle}

O TCE-MA pode utilizar diferentes instrumentos, conforme a finalidade e o regime normativo.

\subsection{Auditoria}

Auditoria é exame sistemático de determinado objeto, orientado por critérios e evidências. Pode incidir sobre conformidade, desempenho, finanças, tecnologia da informação ou outros temas.

\subsection{Inspeção}

Inspeção tende a possuir escopo mais delimitado, destinada a esclarecer fato, suprir omissão ou verificar situação específica.

\subsection{Levantamento}

Levantamento permite conhecer organização, funcionamento, sistemas, programas ou objetos relevantes, identificar riscos e subsidiar planejamento de futuras ações de controle.

\subsection{Acompanhamento}

Acompanhamento observa atividade ou processo enquanto está em execução, permitindo identificar riscos ou irregularidades de forma concomitante.

\subsection{Monitoramento}

Monitoramento verifica o cumprimento de deliberações e a implementação de medidas decorrentes de fiscalizações anteriores.

Os nomes e detalhes procedimentais devem ser compreendidos conforme a Lei Orgânica e o Regimento vigentes, mas a distinção funcional entre conhecer, examinar, acompanhar e verificar cumprimento é essencial.

\section{Organização interna e Regimento Interno}

O Regimento Interno detalha a distribuição de competências entre órgãos colegiados, Presidência, relatoria, unidades técnicas e demais estruturas do Tribunal.

\subsection{Plenário e órgãos colegiados}

As matérias são submetidas ao órgão competente conforme natureza do processo e regras regimentais. Não se deve presumir que todo processo seja decidido pelo Plenário.

\subsection{Relator}

O relator conduz a marcha processual dentro das competências regimentais, examina requerimentos, pode adotar providências instrutórias e apresenta a matéria ao colegiado competente quando necessário.

\subsection{Unidades técnicas}

As unidades técnicas produzem instruções, análises e propostas de encaminhamento. O relatório técnico possui grande importância probatória e argumentativa, mas não se confunde com a decisão do Tribunal.

\subsection{Ministério Público de Contas}

O Ministério Público de Contas atua perante o Tribunal com funções próprias previstas no ordenamento. Não integra a unidade técnica de auditoria e não deve ser confundido com o Ministério Público estadual comum.

\section{Processo de controle externo no TCE-MA}

Embora cada espécie processual possua regras próprias, é útil compreender a lógica geral:
\[
\text{autuação}\rightarrow
\text{distribuição}\rightarrow
\text{instrução}\rightarrow
\text{contraditório}\rightarrow
\text{deliberação}\rightarrow
\text{recurso/execução}\rightarrow
\text{monitoramento}.
\]

\subsection{Instrução}

A instrução reúne documentos, informações de sistemas, análises técnicas, manifestações dos responsáveis e demais evidências necessárias ao julgamento.

\subsection{Contraditório e ampla defesa}

Quando a decisão pode impor débito, multa ou outra consequência desfavorável, deve ser assegurada oportunidade de manifestação ao interessado, sem prejuízo das hipóteses legalmente admitidas de medida cautelar com contraditório diferido.

\subsection{Citação e audiência}

A terminologia e os efeitos devem ser lidos no regime local. Conceitualmente, a citação está associada à formação de relação processual para defesa em situações que podem envolver imputação de débito ou responsabilização; a audiência é usualmente empregada para apresentação de razões sobre irregularidade sem débito. A redação vigente da norma é determinante.

\section{Decisões e efeitos jurídicos}

As deliberações do Tribunal podem assumir efeitos distintos.

\subsection{Contas regulares, regulares com ressalva e irregulares}

A classificação depende da natureza e relevância das ocorrências constatadas. Uma falha formal de pequena materialidade não conduz automaticamente ao mesmo resultado de desfalque, desvio ou dano relevante.

\subsection{Imputação de débito}

Débito está associado a dano ao erário e exige demonstração do prejuízo, quantificação e responsabilização do agente. Não é sinônimo de multa.

\subsection{Multa}

Multa é sanção administrativa prevista em lei. Sua aplicação exige hipótese normativa, competência, contraditório e fundamentação.

\subsection{Determinação e recomendação}

Determinação impõe providência necessária ao cumprimento de obrigação jurídica. Recomendação indica melhoria quando permanece espaço legítimo para escolha administrativa.

\subsection{Medida cautelar}

Cautelar é providência instrumental destinada a evitar dano, agravamento do risco ou perda de utilidade da decisão final. Não constitui sanção antecipada.

\section{Recursos}

Lei Orgânica e Regimento disciplinam espécies recursais, legitimidade, prazos, efeitos e órgão competente. É importante estudar cada recurso como uma estrutura completa:
\[
\text{cabimento} + \text{legitimidade} + \text{prazo} + \text{efeito} + \text{competência}.
\]

O erro comum é memorizar um prazo isolado e transferi-lo para espécie recursal diferente.

\section{Lei Estadual nº 9.936/2013 — organização administrativa}

A Lei nº 9.936, de 22 de outubro de 2013, disciplina a organização administrativa do Tribunal de Contas do Estado do Maranhão, incluindo sua estrutura, cargos e funções administrativas.

Esse diploma não deve ser estudado pela redação original de 2013 de forma isolada. Ele recebeu diversas alterações posteriores e foi novamente modificado em 2026.

\subsection{Alteração pela Lei nº 12.822/2026}

A Lei Estadual nº 12.822, de 30 de março de 2026, alterou a Lei nº 9.936/2013 e também dispositivos relacionados à estrutura do Tribunal. Entre as mudanças expressamente introduzidas estão:
\begin{itemize}
\item acréscimo do \textbf{Gabinete do Procurador-Geral de Contas} à estrutura prevista na Lei nº 9.936/2013;
\item alteração do art. 21, com atualização de valores, quantitativos e regras relativas à Gratificação de Apoio ao Controle Externo;
\item alteração do art. 26 quanto à designação de Auditor Estadual de Controle Externo e Técnico Estadual de Controle Externo para atuação em gabinetes previstos na norma;
\item alteração de anexos que tratam de cargos em comissão e funções gratificadas da estrutura administrativa.
\end{itemize}

\begin{atencao}[redação vigente em 2026]
Para o concurso de 2026, estudar apenas uma versão antiga da Lei nº 9.936/2013 cria risco concreto de erro. Os dispositivos alterados pela Lei nº 12.822/2026 devem ser lidos em sua redação atualizada.
\end{atencao}

\section{Instrução Normativa TCE-MA nº 50/2017 — Tomada de Contas Especial}

A IN nº 50/2017 disciplina medidas administrativas destinadas à elisão de dano e o regime de instauração, pressupostos de constituição, quantificação do débito, conclusão e encaminhamento de tomada de contas especial ao TCE-MA, além de disciplinar o instituto da decadência no contexto tratado pela norma.

\subsection{Finalidade da Tomada de Contas Especial}

A Tomada de Contas Especial (TCE) é instrumento excepcional de apuração de responsabilidade quando existe dano ao erário ou situação legalmente equiparada e as providências administrativas ordinárias não foram suficientes para solucionar a ocorrência.

Sua lógica pode ser representada por:
\[
\text{indício de dano}\rightarrow
\text{medidas administrativas}\rightarrow
\text{persistência da situação}\rightarrow
\text{instauração da TCE}\rightarrow
\text{apuração dos fatos}\rightarrow
\text{responsáveis}\rightarrow
\text{quantificação do débito}\rightarrow
\text{encaminhamento/julgamento}.
\]

A TCE não deve ser utilizada como substituto automático de qualquer procedimento de apuração. A norma privilegia inicialmente providências administrativas capazes de eliminar o dano ou obter ressarcimento.

\subsection{Pressupostos}

A instauração exige identificar o fato gerador e elementos mínimos que indiquem ocorrência de dano, omissão no dever de prestar contas ou outras hipóteses previstas. A mera suspeita genérica, sem delimitação do evento e dos possíveis responsáveis, não constitui instrução adequada.

\subsection{Quantificação do débito}

O processo deve demonstrar como o valor foi apurado, com memória de cálculo, documentos e critérios que permitam reprodução da quantificação. Débito estimado sem base verificável reduz a qualidade da responsabilização.

\subsection{Alterações posteriores}

A IN nº 50/2017 foi alterada por atos posteriores, incluindo a Decisão Normativa nº 28/2017 e a IN nº 56/2018. O sistema eletrônico e-TCEspecial também foi regulamentado para comunicações e encaminhamentos relacionados à tomada de contas especial.

\section{Prescrição, decadência e segurança jurídica}

O controle externo maranhense possui regras próprias sobre decadência e, mais recentemente, disciplina específica sobre prescrição das pretensões punitiva e de ressarcimento. A Resolução TCE-MA nº 383/2023 regulamenta o tema no âmbito do Tribunal.

Prescrição e decadência não são conceitos intercambiáveis. A incidência depende do direito ou pretensão analisada, do termo inicial e das causas de suspensão ou interrupção previstas.

\section{Instrução Normativa TCE-MA nº 82/2025 — emendas parlamentares}

A IN nº 82/2025 disciplina a fiscalização, o acompanhamento e o julgamento da execução de emendas parlamentares estaduais e municipais, com foco em transparência, rastreabilidade e conformidade constitucional.

A norma está inserida em um contexto nacional de reforço da transparência das emendas parlamentares, inclusive em razão das decisões do Supremo Tribunal Federal na ADPF 854.

\subsection{Transparência}

A execução da emenda deve permitir identificar de modo claro a origem dos recursos, autoria da indicação, beneficiário, objeto, valor, execução orçamentária e financeira e resultados associados.

Transparência não é satisfeita apenas pela publicação de valor agregado. É necessária informação suficientemente detalhada para reconstruir o percurso do recurso.

\subsection{Rastreabilidade}

Rastreabilidade significa capacidade de acompanhar o recurso desde a origem até sua aplicação final. Em termos de controle:
\[
\text{autor/origem}\rightarrow
\text{beneficiário}\rightarrow
\text{plano/objeto}\rightarrow
\text{empenho}\rightarrow
\text{pagamento}\rightarrow
\text{execução}\rightarrow
\text{resultado}.
\]

\subsection{Planejamento e plano de trabalho}

A execução deve ser compatível com planejamento e com informações capazes de demonstrar finalidade, objeto e aplicação dos recursos. A existência de transferência financeira não elimina o dever de documentar a destinação e permitir fiscalização.

\subsection{Transferências especiais}

As chamadas \emph{emendas PIX} não constituem recursos livres de controle. A modalidade de transferência pode simplificar determinadas etapas entre entes, mas permanece submetida a princípios constitucionais, transparência, rastreabilidade e fiscalização do uso dos recursos.

\subsection{Sistemas e relatórios eletrônicos}

A lógica da IN exige que informações de execução sejam disponibilizadas de forma apta a permitir acompanhamento e cruzamento pelo Tribunal. Para Auditor/TI, esse ponto aproxima legislação específica de governança de dados: registros incompletos, ausência de identificadores e falta de integração entre execução orçamentária e objeto físico reduzem a rastreabilidade.

\begin{exemplo}[rastreabilidade insuficiente]
Se o portal informa apenas que um município recebeu R$ 1 milhão em transferências especiais, mas não vincula o valor à emenda, ao autor, ao plano de aplicação, aos pagamentos e aos objetos executados, existe transparência nominal, porém rastreabilidade insuficiente para controle substantivo.
\end{exemplo}

\section{Integração entre as normas}

A legislação específica não deve ser estudada como blocos independentes.

Considere uma irregularidade em recurso recebido por emenda parlamentar:
\begin{enumerate}
\item Constituição define princípios e competência de fiscalização;
\item Lei Orgânica define jurisdição, processo e poderes do TCE-MA;
\item Regimento define procedimento interno e competência dos órgãos;
\item IN nº 82/2025 disciplina requisitos específicos de emendas;
\item se houver dano não solucionado administrativamente, o regime de Tomada de Contas Especial pode tornar-se relevante conforme a IN nº 50/2017;
\item eventual responsabilização deve observar contraditório, prescrição e demais normas aplicáveis.
\end{enumerate}

Essa integração é mais importante para resolver casos concretos do que memorizar dispositivos sem compreender sua função.

\section{Síntese conceitual}

\mapafuncional{Legislação do TCE-MA}
{Constituição + Lei Orgânica + Regimento + normas específicas formam um sistema único de controle}
{Estrutura institucional $\rightarrow$ Constituição e Lei Orgânica.\\Organização administrativa $\rightarrow$ Lei 9.936/2013 atualizada pela Lei 12.822/2026.\\Processo e julgamento $\rightarrow$ Lei Orgânica + Regimento.\\Dano/TCE $\rightarrow$ IN 50/2017 e alterações.\\Emendas parlamentares $\rightarrow$ IN 82/2025.}
{Conta de governo: parecer prévio.\\Responsável por gestão: julgamento.\\Dano persistente: medidas administrativas antes da TCE.\\Emenda: exija transparência + rastreabilidade + vinculação ao objeto.\\Norma de 2013: confira alterações de 2026.}
{Auxílio ao Legislativo $\neq$ subordinação.\\Multa $\neq$ débito.\\Cautelar $\neq$ sanção.\\Transferência especial $\neq$ recurso sem controle.\\Lei original $\neq$ redação vigente.}
{Qual norma rege o caso?\\Quem é competente?\\Qual é o tipo de processo?\\Há dano quantificado?\\Existe trilha completa do recurso?\\A redação normativa está atualizada?}

\begin{resumo}
A legislação específica do TCE-MA deve ser dominada como arquitetura jurídica e processual. A Lei Orgânica explica a competência; o Regimento organiza sua execução; a Lei nº 9.936/2013 define estrutura administrativa e deve ser estudada com as alterações de 2026; a IN nº 50/2017 disciplina Tomada de Contas Especial; e a IN nº 82/2025 acrescenta um eixo contemporâneo de fiscalização das emendas parlamentares baseado em transparência e rastreabilidade.
\end{resumo}
